\documentclass[11pt]{article}

% ---------- Packages ----------
\usepackage[a4paper,margin=1in]{geometry}
\usepackage{amsmath,amssymb,amsthm,mathtools}
\usepackage{bbm}
\usepackage{hyperref}
\usepackage{enumitem}
\usepackage{xcolor}
\usepackage{booktabs}

% ---------- Theorem environments ----------
\newtheorem{theorem}{Theorem}
\newtheorem{proposition}{Proposition}
\newtheorem{lemma}{Lemma}
\newtheorem{corollary}{Corollary}

\theoremstyle{definition}
\newtheorem{definition}{Definition}
\newtheorem{remark}{Remark}
\newtheorem{example}{Example}

% ---------- Commands ----------
\newcommand{\Qcal}{\mathcal{Q}}
\newcommand{\Pcal}{\mathcal{P}}
\newcommand{\Bcal}{\mathcal{B}}
\newcommand{\Ccal}{\mathcal{C}}
\newcommand{\Ecal}{\mathcal{E}}
\newcommand{\Fcal}{\mathcal{F}}

\title{
Finitely Additive Observable Laws and the\\
Prokhorov Extension for Query Systems
}

\author{Patrick S. Mahon}

\date{}

\begin{document}

\maketitle

\begin{abstract}
The observational foundations program constructs a canonical probability
space from a compatible family of observable laws, taking countable
additivity of those laws as a hypothesis.  This paper asks when that
hypothesis can itself be derived.  We show that for a countable
Hausdorff query system with continuous refinement maps, any finitely
compatible family of finitely additive normalized valuations satisfying
a \emph{surjective projective uniform tightness} condition admits a
unique extension to a $\sigma$-additive probability measure on the
canonical realization space.  Countable additivity is not assumed
locally --- it is forced globally by the combination of projective
compactness and sequential refinement structure.

The result is a Prokhorov-type extension theorem in the language of
query systems.  It complements the main extension theorem of
\emph{Observational Foundations of Probability} (which takes
$\sigma$-additive marginals as input) by identifying conditions under
which those marginals need not be assumed: they emerge from a tightness
condition on the compatible family as a whole.  The surjective
projective uniform tightness condition is the query-system analogue of
Prokhorov tightness, with the surjectivity of bonding maps on compact
cores playing the role that single-space compactness plays in the
classical theory.
\end{abstract}

%%%%%%%%%%%%%%%%%%%%%%%%%%%%%%%%%%%%%%%%%%%%%%%%%%%%%%%%%%%%%%%%%%%%%
\section{Introduction}
%%%%%%%%%%%%%%%%%%%%%%%%%%%%%%%%%%%%%%%%%%%%%%%%%%%%%%%%%%%%%%%%%%%%%

\paragraph{The foundational question.}
The observational program constructs probability from observation.
Its central extension theorem (\emph{Observational Foundations of
Probability}, hereafter Paper~1) shows that a compatible family of
observable laws $\{\nu_Q\}$ on a query system $(\Qcal, \preceq)$
determines a unique probability measure $P$ on the canonical
realization space $\Omega$.  The construction is a representation
theorem: given that each query $Q$ carries a countably additive law
$\nu_Q$, the global $P$ is assembled from those laws.

The question this paper asks is foundationally prior:

\begin{quote}
\emph{Must one assume countable additivity at each query level, or
can it be derived from the structure of the compatible family?}
\end{quote}

This is Stopping Point~2 of the program: the assumption that each
$\nu_Q$ is $\sigma$-additive is the fuel on which the extension
theorem runs.  The present paper identifies conditions under which
that fuel can itself be derived --- conditions that are properties
of the \emph{system} rather than of individual marginals.

\paragraph{The answer.}
We show that under a condition called \emph{surjective projective
uniform tightness} (SPUT), starting from only finitely additive
normalized valuations $\{\ell_Q\}$ satisfying finite compatibility,
one can derive a $\sigma$-additive global extension.  The argument has
two steps.

The first step is analytic: SPUT supplies compatible compact cores
$\{K_Q\}$ whose inverse limit $K_\Omega \subseteq \Omega$ is compact
and nonempty.  Any decreasing sequence of cylinder sets with uniformly
positive mass must meet $K_\Omega$ at each level; since
$K_\Omega$ is compact, the intersection is nonempty, contradicting
the assumption that the sequence decreases to $\emptyset$.  This
forces the cylinder premeasure $\mu_0$ to be $\sigma$-subadditive,
hence $\sigma$-additive by finite additivity and Carath\'eodory
extension.

The second step is geometric: the extended measure lives on $\Omega$
rather than the ambient product space.  A graph-support lemma shows
that for each comparable pair $Q_1 \preceq Q_2$, the joint law of
$(\mathrm{eval}_{Q_1}, \mathrm{eval}_{Q_2})$ under $P$ is
supported on the graph of $\pi^{Q_2}_{Q_1}$.  For countable
$\Qcal$, a union over countably many null defect sets gives
$P(\Omega) = 1$.

\paragraph{What SPUT is.}
SPUT asks that for every $\varepsilon > 0$ one can find compact sets
$K_Q \subset O_Q$ at each query level such that: (i) the
bonding maps restrict to surjections between compact sets,
$\pi^{Q_2}_{Q_1}(K_{Q_2}) = K_{Q_1}$; and (ii) the mass outside
the compact core is uniformly small, $\ell_Q(O_Q \setminus K_Q) <
\varepsilon$ for every $Q$.  Condition~(i) --- surjectivity of the
bonding maps on the compact cores --- is the critical one.  It ensures
the compact inverse limit $K_\Omega = \varprojlim K_Q$ is nonempty
with surjective projections, which is exactly what the compactness
argument requires.  Without surjectivity, $K_\Omega$ can be empty
even when each $K_Q$ is nonempty.

\paragraph{Comparison to classical Prokhorov.}
In classical Prokhorov tightness~\cite{Prokhorov1956}, one controls mass concentration in
a single space: for every $\varepsilon > 0$, compact $K \subset X$
with $\mu_n(X \setminus K) < \varepsilon$.  The query-system
analogue must control concentration \emph{projectively}, across all
query levels simultaneously, and must ensure the compact sets
assemble coherently in the inverse limit.  SPUT is exactly this
coherent projective tightness condition, with surjectivity of
bonding maps playing the role that single-space compactness plays
classically~\cite{Parthasarathy1967}.

\paragraph{Relation to Paper~1.}
The two extension theorems are complementary:

\begin{center}
\begin{tabular}{lll}
\toprule
& Paper~1 (realizability route) & This paper (tightness route) \\
\midrule
Input measures & $\sigma$-additive $\nu_Q$ & finitely additive $\ell_Q$ \\
Key condition & realizability of $\Omega$ & surjective projective uniform tightness \\
Topology required & no & Hausdorff $O_Q$, continuous $\pi$ \\
Conclusion & $\sigma$-additive $P$ on $\Omega$ & $\sigma$-additive $P$ on $\Omega$ \\
\bottomrule
\end{tabular}
\end{center}

\noindent
Paper~1 takes $\sigma$-additivity as given and derives a global
measure from realizability.  This paper takes only finite additivity
and derives both $\sigma$-additivity and realizability from tightness.
The price is topological structure on the query system.

\paragraph{Structure of the paper.}
Section~\ref{sec:setup} fixes notation and introduces finitely
additive compatible families on query systems.
Section~\ref{sec:sput} defines SPUT, motivates the surjectivity
condition, and gives examples.
Section~\ref{sec:closed} proves $\sigma$-subadditivity on closed
cylinder sequences (the intermediate theorem).
Section~\ref{sec:borel} extends to general Borel cylinders using
Alexandrov regularity.
Section~\ref{sec:concentration} establishes concentration on $\Omega$
via the graph-support lemma.
Section~\ref{sec:main} assembles the full theorem and compares it to
Paper~1.
Section~\ref{sec:discussion} discusses what the result achieves for
the foundational program, what it does not yet achieve, and the
open frontier.

%%%%%%%%%%%%%%%%%%%%%%%%%%%%%%%%%%%%%%%%%%%%%%%%%%%%%%%%%%%%%%%%%%%%%
\section{Setup}\label{sec:setup}
%%%%%%%%%%%%%%%%%%%%%%%%%%%%%%%%%%%%%%%%%%%%%%%%%%%%%%%%%%%%%%%%%%%%%

\begin{definition}[Query system]
A \emph{query system} is a tuple $(\Qcal, \preceq, \{O_Q\},
\{\pi^{Q_2}_{Q_1}\})$ where:
\begin{itemize}
\item $(\Qcal, \preceq)$ is a countable directed preorder (the
  \emph{query index set});
\item each $Q \in \Qcal$ is assigned a measurable outcome space
  $(O_Q, \Bcal(O_Q))$;
\item for each $Q_1 \preceq Q_2$, there is a measurable map
  $\pi^{Q_2}_{Q_1} : O_{Q_2} \to O_{Q_1}$ (the
  \emph{refinement map}) satisfying $\pi^{Q_1}_{Q_1} = \mathrm{id}$
  and $\pi^{Q_3}_{Q_1} = \pi^{Q_2}_{Q_1} \circ \pi^{Q_3}_{Q_2}$
  whenever $Q_1 \preceq Q_2 \preceq Q_3$.
\end{itemize}
The system is \emph{sequentially upper-directed} if for every
countable family $\{Q_n\}_{n \ge 1} \subset \Qcal$ there exists
$Q_* \in \Qcal$ with $Q_* \succeq Q_n$ for all $n$.
\end{definition}

\begin{remark}[Topological enrichment]
Throughout this paper the outcome spaces $O_Q$ are assumed to be
Hausdorff topological spaces with $\Bcal(O_Q)$ the Borel
$\sigma$-algebra, and the refinement maps $\pi^{Q_2}_{Q_1}$ are
assumed to be continuous.  These topological conditions are not
required by Paper~1 and are the price paid for starting from only
finitely additive marginals.
\end{remark}

\begin{definition}[Canonical realization space]
The \emph{canonical realization space} is the projective limit
\[
\Omega
:=
\varprojlim_{Q \in \Qcal} O_Q
=
\left\{
\omega \in \prod_{Q \in \Qcal} O_Q
\;\middle|\;
\pi^{Q_2}_{Q_1}(\omega_{Q_2}) = \omega_{Q_1}
\;\forall\, Q_1 \preceq Q_2
\right\},
\]
equipped with evaluation maps $\mathrm{eval}_Q : \Omega \to O_Q$,
$\omega \mapsto \omega_Q$.

Since each $\pi^{Q_2}_{Q_1}$ is continuous and each $O_Q$ is
Hausdorff, each coherence constraint set
$\{\omega : \pi^{Q_2}_{Q_1}(\omega_{Q_2}) = \omega_{Q_1}\}$ is
closed in the product.  Since $\Qcal$ is countable, $\Omega$ is a
countable intersection of closed sets, hence closed in
$\prod_Q O_Q$.

The \emph{cylinder algebra} $\Ccal$ is generated by sets
$\mathrm{Cyl}(Q, A) := \mathrm{eval}_Q^{-1}(A)$ for $Q \in \Qcal$,
$A \in \Bcal(O_Q)$, and the \emph{observable $\sigma$-field} is
$\sigma(\Qcal) := \sigma(\Ccal)$.
\end{definition}

\begin{definition}[Finitely compatible family]
A family $\{\ell_Q\}_{Q \in \Qcal}$ of normalized valuations
$\ell_Q : \Bcal(O_Q) \to [0,1]$ (with $\ell_Q(O_Q) = 1$) is
\emph{finitely compatible} if for every $Q_1 \preceq Q_2$,
\[
(\pi^{Q_2}_{Q_1})_\# \ell_{Q_2} = \ell_{Q_1},
\]
i.e.\ $\ell_{Q_1}(A) = \ell_{Q_2}((\pi^{Q_2}_{Q_1})^{-1}(A))$ for
all $A \in \Bcal(O_{Q_1})$.  Each $\ell_Q$ is assumed to be
finitely additive (but not necessarily $\sigma$-additive).
\end{definition}

\begin{remark}[The cylinder premeasure]
Given a finitely compatible family $\{\ell_Q\}$, define the
\emph{cylinder premeasure} $\mu_0 : \Ccal \to [0,1]$ by
\[
\mu_0(\mathrm{Cyl}(Q, A)) := \ell_Q(A).
\]
Finite compatibility ensures $\mu_0$ is well-defined (independent of
which query level is used to represent a given cylinder) and
finitely additive on $\Ccal$.  The central question is when $\mu_0$
is $\sigma$-additive, so that Carath\'eodory's theorem~\cite{Billingsley1995,Bogachev2007} extends it
to a probability measure on $\sigma(\Qcal)$.
\end{remark}

%%%%%%%%%%%%%%%%%%%%%%%%%%%%%%%%%%%%%%%%%%%%%%%%%%%%%%%%%%%%%%%%%%%%%
\section{Surjective projective uniform tightness}\label{sec:sput}
%%%%%%%%%%%%%%%%%%%%%%%%%%%%%%%%%%%%%%%%%%%%%%%%%%%%%%%%%%%%%%%%%%%%%

\begin{definition}[Surjective projective uniform tightness]
\label{def:sput}
A finitely compatible family $\{\ell_Q\}$ is \emph{surjectively
projectively uniformly tight} (SPUT) if for every $\varepsilon > 0$
there exists a system $\{K_Q\}_{Q \in \Qcal}$ of compact sets
$K_Q \subset O_Q$ such that:
\begin{enumerate}[label=(\roman*)]
\item \emph{Surjective bonding:}
  $\pi^{Q_2}_{Q_1}(K_{Q_2}) = K_{Q_1}$ for all $Q_1 \preceq Q_2$;
\item \emph{Uniform mass control:}
  $\ell_Q(O_Q \setminus K_Q) < \varepsilon$ for every $Q \in \Qcal$.
\end{enumerate}
\end{definition}

The two conditions play distinct roles.  Condition~(ii) ensures that
the compact cores capture most of the mass at every query level.
Condition~(i) ensures the compact cores assemble coherently across
levels, so their inverse limit is nonempty.

\begin{remark}[Why surjectivity is necessary]
If condition~(i) required only $\pi^{Q_2}_{Q_1}(K_{Q_2}) \subseteq
K_{Q_1}$ (containment rather than equality), the inverse limit
$K_\Omega := \Omega \cap \prod_Q K_Q$ could be empty even when each
$K_Q$ is nonempty.  The classical example: let $O_Q = [0,1]$,
$K_Q = [0, 1-1/Q]$, with refinement maps the identity.  Then each
$K_Q$ is nonempty and the bonding maps satisfy containment, but
$\bigcap_Q K_Q = \emptyset$.  Surjectivity --- $\pi(K_{Q_2})
= K_{Q_1}$ exactly --- rules out such ``shrinking core'' phenomena and
guarantees $K_\Omega \neq \emptyset$ with surjective projections
(Proposition~\ref{prop:compact-core}).
\end{remark}

\begin{remark}[Comparison to classical Prokhorov tightness]
Classical tightness for a family of measures $\{\mu_n\}$ on a single
metric space $X$ asks: for every $\varepsilon > 0$, a compact
$K \subset X$ with $\mu_n(K) \geq 1-\varepsilon$ for all $n$.  SPUT
is the projective analogue: compact sets must be chosen at every
query level simultaneously, and they must assemble into a compact
subset of $\Omega$ via the surjectivity condition.  The surjectivity
of bonding maps is the multi-space analogue of the fact that $K$ is
a single space (automatically surjective over itself) in the
classical case.
\end{remark}

\begin{proposition}[Compact outcome spaces imply SPUT]
\label{prop:compact-implies-sput}
Suppose each outcome space $O_Q$ is compact and each refinement map
$\pi^{Q_2}_{Q_1}$ is surjective.  Then every finitely compatible
family $\{\ell_Q\}$ of normalized valuations is SPUT.
\end{proposition}

\begin{proof}
For any $\varepsilon > 0$, take $K_Q := O_Q$ for every $Q \in \Qcal$.
Then:
\begin{enumerate}[label=(\roman*)]
\item $\pi^{Q_2}_{Q_1}(K_{Q_2}) = \pi^{Q_2}_{Q_1}(O_{Q_2}) = O_{Q_1}
  = K_{Q_1}$ for all $Q_1 \preceq Q_2$, by surjectivity of the
  refinement maps.
\item $\ell_Q(O_Q \setminus K_Q) = \ell_Q(\emptyset) = 0 < \varepsilon$
  for every $Q$.
\end{enumerate}
Both SPUT conditions are satisfied with the bound being $0$, not
merely $< \varepsilon$.
\end{proof}

\begin{remark}[What this means]
When the outcome spaces are compact, the SPUT condition is not a
constraint on the valuations $\{\ell_Q\}$ at all — it is a property
of the query system architecture.  Any finitely compatible family,
however pathological, satisfies SPUT because the compact cores are
simply the full spaces.  The surjectivity of refinement maps then
supplies condition~(i) for free.  No additional tightness assumption
on the observer's laws is required.
\end{remark}

\begin{remark}[Surjective maps without compact spaces]
When the outcome spaces are not compact, surjectivity of the
refinement maps alone does not imply SPUT.  But it does simplify the
construction: given classical tightness of each $\ell_Q$ (compact
$K_Q \subset O_Q$ with $\ell_Q(K_Q) \geq 1 - \varepsilon$), one
can build compatible compact cores by setting $K'_Q :=
\pi^{Q_*}_Q(K_{Q_*})$ for a sufficiently fine $Q_*$.  Surjectivity
ensures $\pi^{Q_2}_{Q_1}(K'_{Q_2}) = K'_{Q_1}$, and the mass
outside the core is controlled by the tightness of $\ell_{Q_*}$.
Thus the delay query systems of \emph{Observational Probability from
Delay Queries} (where refinement maps are surjective coordinate
projections) satisfy SPUT whenever the marginals are individually
tight.
\end{remark}

\begin{definition}[Compact core]
Given a SPUT system $\{K_Q\}$, the \emph{compact core} of $\Omega$
is
\[
K_\Omega
:=
\Omega \cap \prod_{Q \in \Qcal} K_Q
=
\left\{
\omega \in \Omega
\;\middle|\;
\omega_Q \in K_Q \;\forall\, Q \in \Qcal
\right\}.
\]
\end{definition}

\begin{proposition}[Properties of the compact core]
\label{prop:compact-core}
Under SPUT:
\begin{enumerate}[label=(\roman*)]
\item $K_\Omega$ is compact.
\item $K_\Omega$ is nonempty, and $\mathrm{eval}_Q(K_\Omega) = K_Q$
  for every $Q \in \Qcal$.
\end{enumerate}
\end{proposition}

\begin{proof}
(i) The product $\prod_Q K_Q$ is compact by Tychonoff's theorem~\cite{Bourbaki1966}.
Since $\Omega$ is closed in $\prod_Q O_Q$ (as each coherence
constraint is a closed condition), $K_\Omega = \Omega \cap \prod_Q
K_Q$ is a closed subset of a compact space, hence compact.

(ii) The system $(K_Q, \pi^{Q_2}_{Q_1}|_{K_{Q_2}})$ is an inverse
system of nonempty compact Hausdorff spaces.  By condition~(i) of
SPUT, the bonding maps $\pi^{Q_2}_{Q_1}|_{K_{Q_2}} : K_{Q_2} \to
K_{Q_1}$ are surjective.  By the \emph{inverse limit theorem}~\cite{Bourbaki1966,Fremlin2003} for
compact Hausdorff spaces with surjective bonding maps over a directed
index set, $K_\Omega = \varprojlim K_Q$ is nonempty and each
projection $\mathrm{eval}_Q : K_\Omega \to K_Q$ is surjective.
\end{proof}

\begin{remark}[The inverse limit theorem]
Proposition~\ref{prop:compact-core}(ii) uses the classical result
that an inverse system of nonempty compact Hausdorff spaces with
surjective bonding maps has a nonempty inverse limit with surjective
projections.  This follows by expressing $K_\Omega$ as an
intersection in the compact product $\prod_Q K_Q$, where each finite
sub-intersection is nonempty by surjectivity, and applying the
finite intersection property.  The directed-set version of this
argument is standard but is not yet available in Mathlib~4, and
constitutes the main remaining gap in the Lean formalization.
\end{remark}

%%%%%%%%%%%%%%%%%%%%%%%%%%%%%%%%%%%%%%%%%%%%%%%%%%%%%%%%%%%%%%%%%%%%%
\section{The intermediate theorem: closed cylinders}\label{sec:closed}
%%%%%%%%%%%%%%%%%%%%%%%%%%%%%%%%%%%%%%%%%%%%%%%%%%%%%%%%%%%%%%%%%%%%%

The key analytic step is to show that the cylinder premeasure
$\mu_0$ is $\sigma$-subadditive on sequences of closed cylinders.

\begin{theorem}[Closed cylinder $\sigma$-subadditivity]
\label{thm:closed-cylinders}
Let $(\Qcal, \preceq)$ be a countable sequentially upper-directed
query system with Hausdorff outcome spaces and continuous refinement
maps.  Let $\{\ell_Q\}$ be a finitely compatible SPUT family.  If
$C_N = \mathrm{Cyl}(Q_N, F_N) \downarrow \emptyset$ is a decreasing
sequence of cylinder sets with each $F_N \subseteq O_{Q_N}$ closed,
then
\[
\mu_0(C_N) \to 0.
\]
\end{theorem}

\begin{proof}
Suppose for contradiction that $\mu_0(C_N) \geq \varepsilon > 0$
for all $N$.

\textbf{Step 1 (Compact core).}
Apply SPUT with parameter $\varepsilon/2$: obtain compact sets
$\{K_Q\}$ with surjective bonding maps and $\ell_Q(O_Q \setminus
K_Q) < \varepsilon/2$ for all $Q$.

\textbf{Step 2 ($K_\Omega$ is compact and nonempty with surjective
projections).}
By Proposition~\ref{prop:compact-core}, $K_\Omega$ is compact and
$\mathrm{eval}_{Q_N}(K_\Omega) = K_{Q_N}$ for every $N$.

\textbf{Step 3 (Each $C_N \cap K_\Omega$ is nonempty).}
Since $\ell_{Q_N}(F_N) = \mu_0(C_N) \geq \varepsilon$ and
$\ell_{Q_N}(K_{Q_N}) \geq 1 - \varepsilon/2$, we have
\[
\ell_{Q_N}(F_N \cap K_{Q_N})
\geq
\varepsilon - \varepsilon/2 = \varepsilon/2 > 0.
\]
In particular $F_N \cap K_{Q_N} \neq \emptyset$.  Since
$\mathrm{eval}_{Q_N} : K_\Omega \to K_{Q_N}$ is surjective, there
exists $\omega \in K_\Omega$ with $\mathrm{eval}_{Q_N}(\omega) \in
F_N \cap K_{Q_N}$, so $\omega \in C_N \cap K_\Omega$.

\textbf{Step 4 (Each $C_N \cap K_\Omega$ is closed in $K_\Omega$).}
$C_N = \mathrm{eval}_{Q_N}^{-1}(F_N)$ where $F_N$ is closed and
$\mathrm{eval}_{Q_N}$ is continuous, so $C_N$ is closed in $\Omega$
and $C_N \cap K_\Omega$ is closed in $K_\Omega$.

\textbf{Step 5 (Contradiction by compactness).}
The sets $\{C_N \cap K_\Omega\}$ form a decreasing sequence of
nonempty closed subsets of the compact space $K_\Omega$.  By the
finite intersection property,
\[
\bigcap_N (C_N \cap K_\Omega) \neq \emptyset.
\]
But $\bigcap_N C_N = \emptyset$ (since $C_N \downarrow \emptyset$),
so $\bigcap_N (C_N \cap K_\Omega) \subseteq \bigcap_N C_N =
\emptyset$.  Contradiction.
\end{proof}

%%%%%%%%%%%%%%%%%%%%%%%%%%%%%%%%%%%%%%%%%%%%%%%%%%%%%%%%%%%%%%%%%%%%%
\section{Extension to Borel cylinders}\label{sec:borel}
%%%%%%%%%%%%%%%%%%%%%%%%%%%%%%%%%%%%%%%%%%%%%%%%%%%%%%%%%%%%%%%%%%%%%

The intermediate theorem covers closed cylinders.  We now extend to
all Borel cylinders.  The key observation is that no new hypotheses
are needed: inner regularity on each compact core is automatic.

\begin{lemma}[Alexandrov regularity is free]
\label{lem:alexandrov}
Let $K$ be a compact Hausdorff space and $\nu$ a finite Borel measure
on $K$.  Then $\nu$ is inner regular: for every Borel $A \subseteq K$
and $\eta > 0$, there exists closed $F \subseteq A$ with
$\nu(A \setminus F) < \eta$.
\end{lemma}

This is Alexandrov's theorem~\cite{Alexandrov1943,Bogachev2007}; it requires only that $K$ is compact
Hausdorff, which follows from SPUT.

\begin{theorem}[Borel cylinder $\sigma$-subadditivity]
\label{thm:borel-cylinders}
Under the hypotheses of Theorem~\ref{thm:closed-cylinders}, if
$C_N = \mathrm{Cyl}(Q_N, A_N) \downarrow \emptyset$ is a decreasing
sequence of Borel cylinder sets, then $\mu_0(C_N) \to 0$.
\end{theorem}

\begin{proof}
Suppose $\mu_0(C_N) \geq \varepsilon > 0$ for all $N$.
Set $\varepsilon' := \varepsilon/4$.

Apply SPUT with parameter $\varepsilon'$, obtaining compact cores
$\{K_Q\}$ with $\ell_Q(O_Q \setminus K_Q) < \varepsilon'$ for all
$Q$.

For each $N$, apply Alexandrov regularity (Lemma~\ref{lem:alexandrov})
to $\ell_{Q_N}|_{K_{Q_N}}$ with tolerance $\varepsilon'/2^N$: find
closed $F_N \subseteq A_N \cap K_{Q_N}$ (closed in $O_{Q_N}$, since
$K_{Q_N}$ is compact in the Hausdorff space $O_{Q_N}$) with
\[
\ell_{Q_N}(A_N \cap K_{Q_N}) - \ell_{Q_N}(F_N) < \varepsilon'/2^N.
\]
Define closed cylinders $D_N := \mathrm{Cyl}(Q_N, F_N)$.  Then
\[
\mu_0(D_N)
= \ell_{Q_N}(F_N)
\geq \ell_{Q_N}(A_N) - \varepsilon' - \varepsilon'/2^N
\geq \varepsilon - 2\varepsilon'
= \varepsilon/2 > 0.
\]

Let $E_N := \bigcap_{k=1}^N D_k$.  Since the system is sequentially
upper-directed, finite intersections of closed cylinders are again
closed cylinders (each pair $(D_i, D_j)$ has a common refinement
$Q_{ij} \succeq Q_i, Q_j$, giving
$D_i \cap D_j = \mathrm{Cyl}(Q_{ij},
(\pi^{Q_{ij}}_{Q_i})^{-1}(F_i) \cap (\pi^{Q_{ij}}_{Q_j})^{-1}(F_j))$,
which is a closed cylinder).  So each $E_N$ is a closed cylinder.

Moreover, since $C_N$ is decreasing and $D_k \subseteq C_k$:
\[
\mu_0(E_N)
\geq \mu_0(C_N) - \sum_{k=1}^N \mu_0(C_k \setminus D_k)
\geq \varepsilon - \sum_{k=1}^\infty 2\varepsilon'/2^k
= \varepsilon - 2\varepsilon'(1 + 1)
= \varepsilon/2 > 0.
\]
Since $E_N \subseteq C_N$ and $C_N \downarrow \emptyset$, we have
$E_N \downarrow \emptyset$ with each $\mu_0(E_N) \geq \varepsilon/2
> 0$, contradicting Theorem~\ref{thm:closed-cylinders}.
\end{proof}

\begin{corollary}[$\sigma$-additivity of $\mu_0$]
Under the above hypotheses, the cylinder premeasure $\mu_0$ is
$\sigma$-additive on $\Ccal$.  By Carath\'eodory's extension theorem,
$\mu_0$ extends uniquely to a $\sigma$-additive probability measure
on $\sigma(\Qcal)$ over the ambient product $\prod_Q O_Q$.
\end{corollary}

%%%%%%%%%%%%%%%%%%%%%%%%%%%%%%%%%%%%%%%%%%%%%%%%%%%%%%%%%%%%%%%%%%%%%
\section{Concentration on $\Omega$}\label{sec:concentration}
%%%%%%%%%%%%%%%%%%%%%%%%%%%%%%%%%%%%%%%%%%%%%%%%%%%%%%%%%%%%%%%%%%%%%

The Carath\'eodory extension lives on the ambient product
$\prod_Q O_Q$.  We must show it concentrates on the projective limit
$\Omega$.

\begin{lemma}[Graph support]
\label{lem:graph-support}
Let $P$ be the $\sigma$-additive extension of $\mu_0$ to
$\prod_Q O_Q$.  For each $Q_1 \preceq Q_2$, the joint law of
$(\mathrm{eval}_{Q_1}, \mathrm{eval}_{Q_2})$ under $P$ is supported
on the graph
\[
G_{Q_1,Q_2}
:=
\bigl\{
(a,b) \in O_{Q_1} \times O_{Q_2}
\;\big|\;
a = \pi^{Q_2}_{Q_1}(b)
\bigr\}.
\]
\end{lemma}

\begin{proof}
It suffices to verify on measurable rectangles, which generate
$\Bcal(O_{Q_1}) \otimes \Bcal(O_{Q_2})$.  For $A \in \Bcal(O_{Q_1})$
and $B \in \Bcal(O_{Q_2})$:
\begin{align*}
P\bigl(
\mathrm{eval}_{Q_1}^{-1}(A) \cap \mathrm{eval}_{Q_2}^{-1}(B)
\bigr)
&= \mu_0\bigl(\mathrm{Cyl}(Q_2, (\pi^{Q_2}_{Q_1})^{-1}(A) \cap B)\bigr) \\
&= \ell_{Q_2}\bigl((\pi^{Q_2}_{Q_1})^{-1}(A) \cap B\bigr),
\end{align*}
which equals $(\mathrm{graph}_\pi)_\# \ell_{Q_2}(A \times B)$
where $\mathrm{graph}_\pi : b \mapsto (\pi^{Q_2}_{Q_1}(b), b)$.
The measure $(\mathrm{graph}_\pi)_\# \ell_{Q_2}$ is entirely
supported on $G_{Q_1,Q_2}$.

The graph $G_{Q_1,Q_2}$ is Borel: it is the preimage of the diagonal
$\Delta_{O_{Q_1}}$ (closed in Hausdorff $O_{Q_1} \times O_{Q_1}$)
under the continuous map $(a,b) \mapsto (a, \pi^{Q_2}_{Q_1}(b))$.
Hence $G_{Q_1,Q_2}$ is closed, so the statement is well-posed.
\end{proof}

\begin{corollary}
\label{cor:defect-null}
For each $Q_1 \preceq Q_2$, defining the defect set
$D_{Q_1,Q_2} := \{\omega : \omega_{Q_1} \neq
\pi^{Q_2}_{Q_1}(\omega_{Q_2})\}$, we have $P(D_{Q_1,Q_2}) = 0$.
\end{corollary}

\begin{theorem}[Concentration on $\Omega$]
\label{thm:concentration}
$P(\Omega) = 1$.
\end{theorem}

\begin{proof}
The complement decomposes as
$\prod_Q O_Q \setminus \Omega = \bigcup_{Q_1 \preceq Q_2}
D_{Q_1,Q_2}$.
Since $\Qcal$ is countable, the preorder has countably many
comparable pairs.  By Corollary~\ref{cor:defect-null}, each
$D_{Q_1,Q_2}$ is a $P$-null set.  A countable union of null sets
is null, so $P(\prod_Q O_Q \setminus \Omega) = 0$.
\end{proof}

%%%%%%%%%%%%%%%%%%%%%%%%%%%%%%%%%%%%%%%%%%%%%%%%%%%%%%%%%%%%%%%%%%%%%
\section{The main theorem}\label{sec:main}
%%%%%%%%%%%%%%%%%%%%%%%%%%%%%%%%%%%%%%%%%%%%%%%%%%%%%%%%%%%%%%%%%%%%%

\begin{theorem}[Prokhorov extension for query systems]
\label{thm:main}
Let $(\Qcal, \preceq)$ be a countable sequentially upper-directed
query system in which each outcome space $O_Q$ is a Hausdorff
topological space, $\Bcal(O_Q)$ is its Borel $\sigma$-algebra, and
each refinement map $\pi^{Q_2}_{Q_1} : O_{Q_2} \to O_{Q_1}$ is
continuous.  Let $\{\ell_Q\}_{Q \in \Qcal}$ be a finitely compatible
family of finitely additive normalized valuations on $\{(O_Q,
\Bcal(O_Q))\}$.

If $\{\ell_Q\}$ is surjectively projectively uniformly tight
(Definition~\ref{def:sput}), then there exists a unique
$\sigma$-additive probability measure $P$ on
$(\Omega, \sigma(\Qcal))$ satisfying
\[
(\mathrm{eval}_Q)_\# P = \ell_Q
\]
for every $Q \in \Qcal$.  In particular, each $\ell_Q$ is
$\sigma$-additive.
\end{theorem}

\begin{proof}
\emph{$\sigma$-additivity of $\mu_0$.}
By Theorem~\ref{thm:borel-cylinders}, $\mu_0$ satisfies continuity
from above at $\emptyset$ on all Borel cylinders.  Since $\mu_0$ is
already finitely additive, it is $\sigma$-additive on $\Ccal$.

\emph{Carath\'eodory extension.}
By Carath\'eodory's theorem, $\mu_0$ extends uniquely to a
$\sigma$-additive probability measure, still called $P$, on
$\sigma(\Ccal)$ over $\prod_Q O_Q$.

\emph{Concentration.}
By Theorem~\ref{thm:concentration}, $P(\Omega) = 1$.  The
restriction $P|_\Omega$ is the required measure on
$(\Omega, \sigma(\Qcal))$.

\emph{Marginal recovery.}
For any $Q \in \Qcal$ and $A \in \Bcal(O_Q)$:
$P(\mathrm{eval}_Q^{-1}(A)) = \mu_0(\mathrm{Cyl}(Q,A)) = \ell_Q(A)$.

\emph{Uniqueness.}
The observable $\sigma$-field $\sigma(\Qcal)$ is generated by the
cylinder algebra $\Ccal$.  Any two $\sigma$-additive probability
measures on $(\Omega, \sigma(\Qcal))$ with the same $\ell_Q$-marginals
agree on $\Ccal$ and hence on $\sigma(\Qcal)$.
\end{proof}

\begin{remark}[Induced $\sigma$-additivity]
The last sentence of Theorem~\ref{thm:main} is the key point for
the foundational program: the $\ell_Q$ were assumed only finitely
additive, but the theorem produces a $\sigma$-additive global
measure whose marginals are exactly the $\ell_Q$.  In particular,
each $\ell_Q = (\mathrm{eval}_Q)_\# P$ is $\sigma$-additive (as
a pushforward of a $\sigma$-additive measure).  Countable additivity
is not an input --- it is an output.
\end{remark}

\begin{corollary}[Architectural extension theorem]
\label{cor:architectural}
Let $(\Qcal, \preceq)$ be a countable sequentially upper-directed
query system in which each outcome space $O_Q$ is compact Hausdorff
and each refinement map $\pi^{Q_2}_{Q_1}$ is continuous and
surjective.  Then every finitely compatible family $\{\ell_Q\}$ of
finitely additive normalized valuations admits a unique
$\sigma$-additive probability measure $P$ on $(\Omega,\sigma(\Qcal))$
with $(\mathrm{eval}_Q)_\# P = \ell_Q$ for every $Q$.
\end{corollary}

\begin{proof}
By Proposition~\ref{prop:compact-implies-sput}, compact outcome
spaces and surjective refinement maps imply SPUT.
Theorem~\ref{thm:main} then applies directly.
\end{proof}

\begin{remark}[Closing Stopping Point 2]
Corollary~\ref{cor:architectural} is the fully structural version
of Stopping Point~2.  Its hypotheses are entirely architectural:
they concern the topology of the outcome spaces and the surjectivity
of the refinement maps, not the valuations $\{\ell_Q\}$.  Any
observer whose query system has compact outcome spaces and surjective
refinement maps is automatically in the setting where
$\sigma$-additivity is derived from the architecture, not assumed.
The observer need not check whether their laws are tight, or
$\sigma$-additive, or anything beyond finite compatibility.

This is the strongest form of the result: \emph{the structure of
the observational interface alone forces countable additivity of the
resulting global measure.}
\end{remark}

\begin{remark}[Hypothesis audit]
Every hypothesis in Theorem~\ref{thm:main} is used, and each
at exactly one point:

\begin{center}
\begin{tabular}{lll}
\toprule
Hypothesis & Where used & Can it be weakened? \\
\midrule
$\Qcal$ countable & Concentration (\S\ref{sec:concentration})
  & Countable determining subsystem suffices \\
Sequentially upper-directed & Borel extension (\S\ref{sec:borel})
  & Needed for closed-cylinder intersections \\
Hausdorff $O_Q$ & $\Omega$ closed; graphs Borel
  & Essential \\
Continuous $\pi$ & $\Omega$ closed; cylinders closed
  & Essential \\
Surjective bonding on $K_Q$ & $K_\Omega \neq \emptyset$
  & Cannot weaken to containment \\
Compact $K_Q$ & Tychonoff; FIP; Alexandrov
  & Essential \\
$\ell_Q(O_Q \setminus K_Q) < \varepsilon$ & Mass estimates
  & Essential \\
Finite compatibility & $\mu_0$ well-defined
  & Essential \\
\bottomrule
\end{tabular}
\end{center}
\end{remark}

%%%%%%%%%%%%%%%%%%%%%%%%%%%%%%%%%%%%%%%%%%%%%%%%%%%%%%%%%%%%%%%%%%%%%
\section{Discussion}\label{sec:discussion}
%%%%%%%%%%%%%%%%%%%%%%%%%%%%%%%%%%%%%%%%%%%%%%%%%%%%%%%%%%%%%%%%%%%%%

\paragraph{What the theorem achieves.}
Theorem~\ref{thm:main} addresses Stopping Point~2 of the
observational program: it shows that $\sigma$-additivity need not
be assumed at each query level.  Instead, it is a global consequence
of two structural conditions: finite compatibility (the marginals
cohere under refinement) and SPUT (the mass concentrates
projectively on compatible compact cores).  These conditions are
properties of the \emph{system} of observable laws, not of
individual marginals.

The result is a strict weakening of the input to Paper~1:
Paper~1 assumes $\sigma$-additive marginals and realizability,
while this paper needs only finitely additive marginals and SPUT.
The conclusion is the same: a $\sigma$-additive global measure on
$(\Omega, \sigma(\Qcal))$.

\paragraph{What SPUT is buying.}
SPUT is a condition about how observable mass concentrates as one
moves between query levels.  It says, in operational terms: the
things the observer considers \emph{likely} at fine resolutions
assemble coherently into things that are likely at coarser
resolutions, and this assembly is non-degenerate (surjectivity).
This is a natural condition on a query system arising from a
physical process: if the system concentrates mass on a compact
region at every observable resolution, that region is observable
at all resolutions jointly.

\paragraph{What it costs.}
SPUT requires topological structure: Hausdorff spaces and continuous
refinement maps.  This is the price for starting from finitely
additive marginals.  Paper~1 operates without this topology, at
the cost of requiring $\sigma$-additivity as an input.  The two
theorems occupy adjacent positions in the space of hypotheses,
with this paper trading topology for weaker input measures.

\paragraph{The deeper question.}
Stopping Point~2 in its strongest form asks whether $\sigma$-additivity
can be forced from \emph{purely combinatorial} conditions on the
query system, without topological assumptions.  Theorem~\ref{thm:main}
does not achieve this: it replaces the assumption of
$\sigma$-additivity with the assumption of Hausdorff topology plus
SPUT.  The question is whether SPUT itself can be derived from
something more primitive.

\paragraph{The architectural theorem and the deeper question.}
Corollary~\ref{cor:architectural} closes the first form of
Stopping Point~2: in the compact outcome space setting with
surjective refinement maps, $\sigma$-additivity is a theorem of
architecture.  The observer contributes only finite compatibility;
the system does the rest.

This raises the next question: how restrictive is the compactness
assumption?  For sensor systems with bounded output spaces (voltages,
pressures, counts in a finite range), compactness is natural.  For
infinite-range observables, it must be imposed or approximated.
Understanding when the compact-outcome-space setting is appropriate,
and what weaker conditions might substitute for it, is the next
foundational target.

\paragraph{Open problems.}

\begin{enumerate}
\item \emph{The uncountable case.}
Concentration (Theorem~\ref{thm:concentration}) uses countability
of $\Qcal$ to take a countable union of null defect sets.  For
uncountable $\Qcal$, this argument fails.  The correct
generalization appears to be a condition on determining subsystems:
if $\Qcal$ admits a countable determining subsystem $\Qcal_0$
(coherence on $\Qcal_0$ forces coherence on all of $\Qcal$), then
the countable argument applies.  Characterizing when such subsystems
exist is the open problem.

\item \emph{Topology-free route.}
Theorem~\ref{thm:main} requires Hausdorff topology and continuity
of refinement maps.  A topology-free version would need to replace
Alexandrov's theorem (automatic inner regularity on compact
Hausdorff spaces) with a purely measure-theoretic substitute.
Whether this is possible --- and what structural condition on the
query system it would require --- remains open.

\item \emph{Lean formalization.}
The main outstanding gap in the Lean~4/Mathlib formalization is
the directed-set form of the inverse limit theorem
(Proposition~\ref{prop:compact-core}(ii)).  Mathlib~4 does not
yet contain this result in the form needed.  All other steps of
the proof have been formalized or have clear paths to
formalization once this gap is closed.
\end{enumerate}

\bibliographystyle{plain}
\bibliography{references}

\section*{Acknowledgements}

The author thanks Claude (Anthropic) for assistance with mathematical
writing and technical discussion throughout the development of this
work.  The author gratefully acknowledges financial support from
Dr.\ Simon Bonner and Dr.\ Matt Davison at the University of Western
Ontario, and from MITACS.

\end{document}
